\section{Introduction and Project Overview}
\label{sec:introduction}

\sectionnote{Owner: Sydney Eichenberg. Target: 0.5 page. Rubric: Good report organization break-down, within page limits.}

\begin{requirements}
  \item ``Great'' on this criterion means a clear structure and logical flow, with no repetition, within the page limit.
  \item A reader of the introduction alone should know what the report found and where each part is.
\end{requirements}

\subsection{Purpose and Scope}
\label{sec:purpose}

\begin{requirements}
  \item State the goal of the report: recover WezTerm's concrete architecture, then identify and explain where it diverges from the conceptual architecture.
  \item State exactly what was analysed: the WezTerm repository at commit \pinnedcommit{} (August~31, 2026) and the course-provided Understand project built from it, both cited.
  \item State what the analysis builds on: our A1 conceptual architecture, cited, as revised in Section~\ref{sec:conceptual}.
  \item State the scope: the whole system at the top level, and one subsystem (Section~\ref{sec:subsystem}) at the second level.
\end{requirements}

\subsection{Report Organization}
\label{sec:organization}

\begin{requirements}
  \item One sentence per section, in the order they appear, ending with the AI appendix.
\end{requirements}

\subsection{Summary of Findings}
\label{sec:summary-findings}

\begin{requirements}
  \item The three or four most important findings: how closely the concrete architecture matches the conceptual one, the most significant divergences and their causes, and the main result for the second-level subsystem.
  \item Keep these consistent with the Abstract and the Conclusions.
\end{requirements}
